\documentclass[12pt]{article}

% The preamble file contains all the macros and page setup. I recommend making a folder where you compile your exercises and always call on the preamble.
%%%%%%%%%%%%%%%%%%%%%%%%%%%%%%%%%%%%%%%%%%%%%%%%%%%%%%%%%%%%%%%%%%%%%%
%% all the formatting stuff and packages



\usepackage{geometry}
\geometry{
  top=4cm,
  headheight=70pt,
  headsep=25pt,
  inner=1.25in,
  outer=1.25in
}

\usepackage[utf8]{inputenc}

\usepackage{xcolor,graphicx,hyperref,amsthm,amsmath,amssymb}
\usepackage{algpseudocode}


\usepackage{fancyhdr}
\pagestyle{fancy}

%%% Begin Header Definitions %%%
\fancypagestyle{firststyle}
{
   	\lhead{}
	\rhead{\raisebox{0.2cm}{\includegraphics[width=3.5cm]{resource/hpi_logo.pdf}}}
	\renewcommand{\headrulewidth}{0pt}
}


\newcommand{\makeheader}[1]{
\noindent
\textbf{Algorithmix \hfill Winter term 2019/20 \\ Andreas G\"obel}

\begin{center}
\large{\textbf{\makeHeadline}}\\
\end{center}

\begin{center}
\textit{#1} 
\end{center}


}

\newcommand{\lecturer}[1]{\newcommand{\makeLecturer}{#1}}
\newcommand{\lectureName}[1]{\newcommand{\makeLectureName}{#1}}
\newcommand{\semester}[1]{\newcommand{\makeSemester}{#1}}
\newcommand{\headline}[1]{\newcommand{\makeHeadline}{#1}}
\newcommand{\urlLine}[1]{\newcommand{\makeUrlLine}{#1}}
%%% End Header Definitions %%%

\newcommand{\pr}{\ensuremath{\mathbb{P}}}
\newcommand{\E}{\ensuremath{\mathbb{E}}}


\newcommand{\makeCompact}{%
\setlength{\itemsep}{1pt}%
\setlength{\parsep}{0pt}%
\setlength{\parskip}{0pt}}


\newcommand{\Var}{\mathrm{Var}}


\newcommand{\student}[2]{
\thispagestyle{firststyle}

%%% Begin Header Specification %%%

\headline{Solutions to homework #1}

\makeheader{#2}

\medskip

}

\newcommand{\sol}[1]{ \hspace{.5cm}\\ \noindent \textbf{Exercise #1.}}

% Here you can define your own macros
\newcommand{\calF}{\ensuremath{\mathcal{F}}}

\begin{document}
 
 %To create a header for the document use the following macro. \student{Homework number}{Name}

 \student{3}{Stefan Reschke}
 

 \sol{1} (a) If $n=3$ with coin values $[1,3,4]$ and try the algorithm on the value $6$, we will get $[4,1,1]$ which is of length $3$ but there is also the shorter $[3,3]$ with a length of $2$. Thus this capitalist algorithm will not give smallest coin combination for all possible coin values.
 
 (b) Given a value $k$ we know that \begin{equation}
 	\exists l\in\mathbb{N}:\ b^l<k\ \land\ b^{l+1}>k.
 \end{equation} Thus by picking $b^l$ we reduce from $k$ and call with $k-b^l$ recursively. In this step, we can either pick $b^l$ again or use a coin of even lower value. This way we efficiently reduce to the space between $1$ and $b$.
 
 If $1\leq k\leq b-1$ we always need to pick the coin of value $1$ exactly $k$ times. There is no more efficient way to give those values. Thus the algorithm always works in this coin set.
 
 (c) Again on a value $k$, we can call recursively on all values on $c_i<k$. The smallest number of coins to give the value $k$ must be \begin{equation}\label{1c:dynprog-rule}
  	\texttt{nbCoins}(k):=\texttt{min}\big\{\texttt{nbCoins}(k-b)\ |\ b\in c\big\} + 1.
 \end{equation} Thus we have a rule that reduces to subproblems. We apply dynamic programming and store the results of the subproblems in an array. We use the algorithm in \autoref{algo:1}.
 
 \textbf{Is it correct?} Since $c[1]=1$, the second loop will always be run and all entries in $a$ will be written with the correct value according to \autoref{1c:dynprog-rule}. So: yes!
 
 \textbf{How fast is it?} It will take at most $\#c\cdot T$ steps, thus $\mathcal{O}(\#c\cdot T)$.
 
 \textbf{Can we do better?} No subproblem is calculated a second time. I think we can't do better.

\begin{lstlisting}[caption={Minimum number of coins},label={algo:1}]
Name: nbCoins
Input: target value T, coins c
Output: number of coins in c to give T

a := Array of length T filled with infinities.
for t in {1, ..., T}:
	for k in c, if k < t:
		new := a[t-k] + 1.
		if new < a[t]:
			a[t] := new.
return a[T]
\end{lstlisting}
 
 \sol{3} This feels a lot like quicksort. If the cuts closest to the center are done first and then again first for each recursive call, we get the optimal result. If our algorithm \texttt{cut} takes a length $l$ of the array at hand and a list of of points $p$ where to cut, we have \begin{multline}
 	\texttt{cut}(l, P)=\texttt{min}_{p\in P}\bigg(\texttt{cut}\big(p, \{k\in P\ |\ k<p\}\big) \\
 	+\texttt{cut}\big(l-p, \{k-p\ |\ k\in P\ \land\ k>p\}\big)\bigg)+l.
 \end{multline}
 
 We can use the algorith in \autoref{algo:3}. It takes a length $l$ and a set $P$ as parameters and a dictionary $T$ to store results of the sub problems. At the start all keys of $T$ are empty. The overall result is stored in $T[l,P]$.

 \vspace{.5cm}
 \textbf{Is it correct?} All we did was applying the above rule. Since its correct, the algorithm also must give correct results.
 
\begin{lstlisting}[caption={Cutting arrays recursively},label={algo:3}]
Name: cut
Input: length of array l, set of cuts P, dictionary T
Output: cost of performing cuts in P on array of length l
 
if P is empty:
	T[l, P] := 0.
	return.
result := infinity.
for p in P:
	Left := {k in P | k < p}.
	Right := {k - p | k in P & k > p}.
	cut(p, Left, T)
	cut(l-p, Right, T)
	result := min(result, T[p, Left] + T[l-p, Right]).
T[l, P] := result + l.
\end{lstlisting}

 \sol{4} Our goal is to maximise $P[x,y]$ for all $x\times y$ blocks. Let $M[x,y]$ the maximum price we can achieve for a (split) block of lengths $x,y$. Each time, we have three choices: \begin{enumerate}
 	\item Leave it as it is (if $M[x,y]=P[x,y]$)
 	\item Cut at $x$, at a line $i$ s.t. $M[i,y]+M[x-i,y]>P[x,y]$
 	\item Cut at $y$, at a line $i$ s.t. $M[x,i]+M[x,y-i]>P[x,y]$
 \end{enumerate}
 Thus we can split it up into smaller sub problems. For the latter cases, we only have to look at $i$s where $1\leq i\leq\floor{\frac{x}{2}}$ or $1\leq i\leq\floor{\frac{y}{2}}$ respectively, because addition is commutative. Thus we have \begin{multline}
 	M[x,y]=\texttt{max}\bigg\{P[x,y],\ \texttt{max}_{1\leq i\leq\floor{\frac{x}{2}}} (M[i,y]+M[x-i,y]),\ \texttt{max}_{1\leq i\leq\floor{\frac{y}{2}}} (M[x,i]+M[x,y-i])\bigg\}.
 \end{multline} This we can again forge into an recursive dynamic programming algorithm. We use the algorithm in \autoref{algo:4} and say that its $x\times y$-table $M$ is initialized with empty values and the result of the overall algorithm is written to $M[x,y]$.

 \vspace{.5cm}
 \textbf{Is it correct?} Recursion stops if we cut to lengths of $0$, thus we still keep cutting plates of length $1$ into smaller plates. Everything we did was formalizing the rule above. Therefore this gives correct results.
 
 \textbf{How much time does it take?} $\mathcal{O}(x\cdot y)$ as we do not call a sub problem twice (see first assertion in algorithm).

 \textbf{Can we do better?} There probably is a nicer way of writing this algorithm. But I think we can't do any better because we cannot make assumptions on $P[x,y]$.

\begin{lstlisting}[caption={Cutting marble recursively},label={algo:4}]
Name: optimalCut
Input: lengths x and y, table M
Output: nothing, fills table M

if M[n,m] already set:
    return.
if x=0  or  y=0 then:
    M[i,j] := 0.
else:
    max := 0.
    for i in {1, ..., floor(x/2)}:
        optimalCut(i, y, M).
        optimalCut(x-i, y, M).
        if M[i,y] + M[x-i,y] > max:
            max := M[i,y] + M[x-i,y].
    for i in {1, ..., floor(y/2)}:
        optimalCut(x, i, M).
        optimalCut(x, y-i, M).
        if M[x,i] + M[x,y-i] > max:
            max := M[x,i] + M[x,y-i].
    if P[x,y] > max:
        max := P[x,y].
    M[n][m] := max.
\end{lstlisting}
 
\end{document}
